\documentclass[11pt]{article}

% ----------------------------------------------------------------------------
%  RSA and Shor's Algorithm — a detailed lesson
%  Companion to IBM/22_Shor_Period_Finding_to_RSA.ipynb (and IQM/22_Shor_Period_Finding_to_RSA.ipynb) in the QC-Course repository.
%  Compile with:  pdflatex rsa_shor_lesson.tex   (twice, for cross-references)
% ----------------------------------------------------------------------------

\usepackage[utf8]{inputenc}
\usepackage[T1]{fontenc}
\usepackage{amsmath,amssymb,amsthm}
\usepackage{mathtools}
\usepackage[margin=1in]{geometry}
\usepackage{enumitem}
\usepackage{xcolor}
\usepackage[most]{tcolorbox}
\usepackage{graphicx}
\usepackage[colorlinks=true,linkcolor=blue!55!black,urlcolor=blue!55!black,citecolor=blue!55!black]{hyperref}

% --- theorem environments ----------------------------------------------------
\theoremstyle{plain}
\newtheorem{theorem}{Theorem}[section]
\newtheorem{lemma}[theorem]{Lemma}
\newtheorem{corollary}[theorem]{Corollary}
\newtheorem{proposition}[theorem]{Proposition}
\theoremstyle{definition}
\newtheorem{definition}[theorem]{Definition}
\newtheorem{example}[theorem]{Example}
\newtheorem{recipe}[theorem]{Recipe}
\theoremstyle{remark}
\newtheorem{remark}[theorem]{Remark}

% --- a boxed recipe ----------------------------------------------------------
\newtcolorbox{recipebox}[1]{colback=gray!4,colframe=black!60,
  fonttitle=\bfseries,title={#1},breakable,boxrule=0.6pt,arc=2pt,left=4pt,right=4pt}

% --- handy macros ------------------------------------------------------------
\newcommand{\Z}{\mathbb{Z}}
\newcommand{\N}{\mathbb{N}}
\newcommand{\Zn}{(\Z/n\Z)^{\!*}}
\newcommand{\gcdf}[2]{\gcd(#1,#2)}
\newcommand{\ket}[1]{\lvert #1 \rangle}
\newcommand{\bra}[1]{\langle #1 \rvert}
\newcommand{\modN}{\ (\mathrm{mod}\ N)}
\DeclareMathOperator{\lcm}{lcm}
\DeclareMathOperator{\ord}{ord}

\title{\textbf{RSA and Shor's Algorithm}\\[2pt]
\large From Number Theory to Quantum Factoring}
\author{QC-Course 2026 --- lesson notes\\[2pt]
\small companion to \texttt{IBM/16\_Shor\_Period\_Finding\_to\_RSA.ipynb} and \texttt{IQM/16\_Shor\_Period\_Finding\_to\_RSA.ipynb}}
\date{\today}

\begin{document}
\maketitle

\begin{abstract}
\noindent These notes explain, from first principles, how \emph{asymmetric}
(public-key) cryptography works, why RSA in particular is secure against classical
computers, and how a quantum computer running \emph{Shor's algorithm} breaks it.
We state and explain every number-theoretic theorem RSA relies on (Fermat, Euler,
B\'ezout, the Chinese Remainder Theorem), prove that RSA decryption inverts
encryption, and then build Shor's algorithm step by step: the reduction from
factoring to order finding, quantum phase estimation of the modular-multiplication
operator, and continued-fraction post-processing. Practical ``recipes'' and fully
worked numerical examples accompany each step, and map directly onto the cells of
the companion Jupyter notebook.
\end{abstract}

\tableofcontents
\bigskip

% ============================================================================
\section{Symmetric vs.\ asymmetric cryptography}
% ============================================================================

In a \emph{symmetric} cipher (AES, ChaCha20, \dots) the same secret key $k$ both
encrypts and decrypts: $c = E_k(m)$, $m = D_k(c)$. This is fast, but it has a famous
weakness --- the \emph{key-distribution problem}: before Alice and Bob can talk
securely, they must already share $k$ over some secure channel. For $n$ mutually
communicating parties one needs $\binom{n}{2}$ shared keys.

\emph{Asymmetric} (public-key) cryptography, introduced by Diffie and Hellman (1976)
and first realised concretely by Rivest, Shamir and Adleman (1977), removes this
problem. Each party publishes a \textbf{public key} $\mathrm{pk}$ and keeps a
\textbf{private key} $\mathrm{sk}$. Anyone can encrypt \emph{to} Bob using his public
key, but only Bob --- holder of the matching private key --- can decrypt:
\[
  c = E_{\mathrm{pk}}(m), \qquad m = D_{\mathrm{sk}}(c).
\]

\begin{definition}[One-way and trapdoor functions]
A function $f$ is \emph{one-way} if $y=f(x)$ is easy to compute but $x=f^{-1}(y)$ is
computationally infeasible. It is a \emph{trapdoor} one-way function if there is extra
secret information (the trapdoor) that makes inverting $f$ easy again.
\end{definition}

Public-key cryptography needs trapdoor one-way functions. For RSA the trapdoor is the
\textbf{factorisation of the modulus}: multiplying two large primes $p,q$ to get
$N=pq$ is easy, but recovering $p,q$ from $N$ is (classically) believed hard --- unless
you know the trapdoor. Shor's algorithm destroys this asymmetry on a quantum computer.

% ============================================================================
\section{Number-theoretic background}
% ============================================================================

All of RSA lives in \emph{modular arithmetic}. We write $a \equiv b \pmod n$ to mean
$n \mid (a-b)$, i.e.\ $a$ and $b$ leave the same remainder on division by $n$. The set
of residues $\Z/n\Z = \{0,1,\dots,n-1\}$ is a commutative ring under $+$ and $\times$
modulo $n$.

\subsection{gcd, B\'ezout, and modular inverses}

\begin{theorem}[B\'ezout's identity]\label{thm:bezout}
For integers $a,b$ not both zero, there exist integers $x,y$ with
\[
  ax + by = \gcd(a,b).
\]
Moreover $x,y$ are computed in $O(\log \min(a,b))$ steps by the
\emph{extended Euclidean algorithm}.
\end{theorem}

\begin{proof}[Idea]
The set $S=\{ax+by : x,y\in\Z\}\cap \N_{>0}$ is non-empty; let $d=ax_0+by_0$ be its
least element. Dividing $a$ by $d$ with remainder, $a=qd+r$ with $0\le r<d$, gives
$r=a-qd=a(1-qx_0)-b(qy_0)\in S\cup\{0\}$. Minimality of $d$ forces $r=0$, so $d\mid a$
and likewise $d\mid b$; hence $d\mid\gcd(a,b)$. Conversely $\gcd(a,b)\mid d$ because it
divides every term $ax+by$. Thus $d=\gcd(a,b)$.
\end{proof}

\begin{corollary}[Existence of modular inverses]\label{cor:inverse}
An element $a\in\Z/n\Z$ has a multiplicative inverse (an $a^{-1}$ with
$a\,a^{-1}\equiv 1 \pmod n$) \textbf{iff} $\gcd(a,n)=1$. The inverse is the coefficient
$x$ in B\'ezout's identity $ax+ny=1$, reduced modulo $n$.
\end{corollary}

This corollary is exactly what lets us turn the public exponent $e$ into the private
exponent $d$ in RSA: $d \equiv e^{-1} \pmod{\varphi(N)}$.

\begin{definition}[Group of units]
The \emph{units} of $\Z/n\Z$ are the invertible residues,
$\Zn = \{\, a : 1\le a\le n,\ \gcd(a,n)=1 \,\}$.
This is an abelian group under multiplication modulo $n$.
\end{definition}

\subsection{Euler's totient function}

\begin{definition}[Euler's totient]
$\varphi(n) = \bigl|\Zn\bigr|$ is the number of integers in $\{1,\dots,n\}$ coprime
to $n$.
\end{definition}

\begin{proposition}[Totient of a product of distinct primes]\label{prop:phi}
If $p\neq q$ are primes then $\varphi(pq) = (p-1)(q-1)$. More generally $\varphi$ is
\emph{multiplicative}: $\varphi(mn)=\varphi(m)\varphi(n)$ whenever $\gcd(m,n)=1$, and
$\varphi(p)=p-1$ for a prime $p$.
\end{proposition}

\begin{proof}[Sketch]
Among $\{1,\dots,pq\}$ the numbers \emph{not} coprime to $pq$ are the $q$ multiples of
$p$ and the $p$ multiples of $q$, with the single number $pq$ counted twice. Hence
$\varphi(pq)=pq-p-q+1=(p-1)(q-1)$. Multiplicativity follows from the Chinese Remainder
Theorem below, which gives a ring isomorphism
$\Z/mn\Z\cong \Z/m\Z\times\Z/n\Z$ and hence a bijection of unit groups.
\end{proof}

\subsection{Fermat's little theorem and Euler's theorem}

\begin{theorem}[Fermat's little theorem]\label{thm:fermat}
If $p$ is prime and $p\nmid a$, then $a^{p-1}\equiv 1 \pmod p$. Equivalently,
$a^{p}\equiv a \pmod p$ for \emph{every} integer $a$.
\end{theorem}

\begin{theorem}[Euler's theorem]\label{thm:euler}
If $\gcd(a,n)=1$ then
\[
  a^{\varphi(n)} \equiv 1 \pmod n .
\]
\end{theorem}

\begin{proof}
Work in the finite abelian group $\Zn$ of order $\varphi(n)$. By Lagrange's theorem the
order of any element divides the order of the group, so $a^{\varphi(n)}=1$ in $\Zn$,
i.e.\ $a^{\varphi(n)}\equiv 1 \pmod n$. (Fermat's little theorem is the special case
$n=p$ prime, where $\varphi(p)=p-1$.) \end{proof}

Euler's theorem is the engine of RSA: raising to the power $\varphi(N)$ acts like the
identity on units, so exponents can be reduced modulo $\varphi(N)$.

\subsection{The Chinese Remainder Theorem}

\begin{theorem}[CRT]\label{thm:crt}
If $\gcd(m,n)=1$, the system $x\equiv a \pmod m,\ x\equiv b \pmod n$ has a unique
solution modulo $mn$. Equivalently, the map
\[
  \Z/mn\Z \;\longrightarrow\; \Z/m\Z \times \Z/n\Z, \qquad
  x \longmapsto (x \bmod m,\ x\bmod n)
\]
is a ring isomorphism.
\end{theorem}

CRT is used both to \emph{prove} RSA correct (Section~\ref{sec:rsa-correct}) and to
\emph{speed it up} in practice (decrypting modulo $p$ and $q$ separately is $\sim 4\times$
faster than modulo $N$).

% ============================================================================
\section{The RSA cryptosystem}
% ============================================================================

\subsection{Key generation, encryption, decryption}

\begin{recipebox}{Recipe R1 --- RSA key generation}
\begin{enumerate}[leftmargin=2.2em,itemsep=1pt,topsep=3pt]
  \item Choose two large distinct primes $p,q$ (e.g.\ 1024 bits each for a 2048-bit key).
  \item Compute the modulus $N = pq$ and the totient $\varphi(N) = (p-1)(q-1)$.
  \item Choose a public exponent $e$ with $1<e<\varphi(N)$ and $\gcd(e,\varphi(N))=1$
        (commonly $e=65537$).
  \item Compute the private exponent $d \equiv e^{-1} \pmod{\varphi(N)}$ via the
        extended Euclidean algorithm (Corollary~\ref{cor:inverse}).
  \item \textbf{Public key} $\mathrm{pk}=(N,e)$. \quad
        \textbf{Private key} $\mathrm{sk}=(N,d)$. \\
        Discard (or guard) $p,q,\varphi(N)$ --- knowing any of them breaks the scheme.
\end{enumerate}
\end{recipebox}

For a message represented as an integer $m$ with $0\le m<N$:
\[
  \text{(encrypt)}\quad c = m^{e} \bmod N,
  \qquad\qquad
  \text{(decrypt)}\quad m = c^{d} \bmod N.
\]

\subsection{Why decryption works}\label{sec:rsa-correct}

\begin{theorem}[RSA correctness]\label{thm:rsa}
With the parameters of Recipe~R1, for every integer $m$ with $0\le m<N$,
\[
  (m^{e})^{d} \equiv m \pmod N .
\]
\end{theorem}

\begin{proof}
By construction $ed\equiv 1 \pmod{\varphi(N)}$, so $ed = 1 + k\varphi(N)$ for some
integer $k\ge 0$, where $\varphi(N)=(p-1)(q-1)$. We show $m^{ed}\equiv m$ modulo each
prime and combine with CRT (Theorem~\ref{thm:crt}).

\emph{Modulo $p$.} If $p\mid m$ then $m^{ed}\equiv 0\equiv m \pmod p$. Otherwise
$\gcd(m,p)=1$, and Fermat's little theorem (Theorem~\ref{thm:fermat}) gives
$m^{p-1}\equiv 1\pmod p$. Since $(p-1)\mid \varphi(N)$,
\[
  m^{ed}=m^{1+k\varphi(N)}
       = m\cdot\bigl(m^{\,p-1}\bigr)^{k(q-1)}
       \equiv m\cdot 1^{\,k(q-1)} = m \pmod p .
\]
\emph{Modulo $q$.} The same argument gives $m^{ed}\equiv m \pmod q$.

Since $p\neq q$ are coprime and $m^{ed}-m$ is divisible by both $p$ and $q$, it is
divisible by $N=pq$; hence $m^{ed}\equiv m \pmod N$.
\end{proof}

\subsection{A fully worked (small) example}

\begin{example}[RSA by hand]\label{ex:rsa}
Take $p=5$, $q=11$. Then
\[
  N = 55, \qquad \varphi(N) = (5-1)(11-1) = 40 .
\]
Choose $e=3$ (indeed $\gcd(3,40)=1$). The private exponent solves $3d\equiv 1 \pmod{40}$;
since $3\cdot 27 = 81 = 2\cdot 40 + 1$, we get $d=27$.
\[
  \mathrm{pk}=(55,3),\qquad \mathrm{sk}=(55,27).
\]
Encrypt $m=7$:
\[
  c = 7^{3}\bmod 55 = 343 \bmod 55 = 13 .
\]
Decrypt $c=13$ (using CRT to keep the numbers small):
\[
  13 \equiv 3 \!\!\pmod 5,\quad 3^{27}\equiv 3^{27 \bmod 4}=3^{3}=27\equiv 2\!\!\pmod 5;
\]
\[
  13 \equiv 2 \!\!\pmod{11},\quad 2^{27}\equiv 2^{27 \bmod 10}=2^{7}=128\equiv 7\!\!\pmod{11}.
\]
CRT reassembles $x\equiv 2 \pmod 5$, $x\equiv 7 \pmod{11}$ into $x=7$. Decryption
returns the original $m=7$. \checkmark
\end{example}

\begin{remark}[The textbook ``big'' example]
The classic RSA example (and the one in \texttt{18\_RSA.ipynb}) uses $p=61$, $q=53$,
so $N=3233$, $\varphi(N)=3120$, $e=17$, $d=2753$. Encrypting $m=65$ gives
$c=65^{17}\bmod 3233 = 2790$, and $2790^{2753}\bmod 3233 = 65$ again.
\end{remark}

\subsection{Security: what an attacker must do}

\begin{proposition}[Knowing $\varphi(N)$ is as hard as factoring]\label{prop:phi-fact}
Given $N=pq$, recovering $\varphi(N)=(p-1)(q-1)$ is computationally equivalent to
factoring $N$.
\end{proposition}

\begin{proof}
If $p,q$ are known, $\varphi(N)$ is immediate. Conversely, from $N=pq$ and
$\varphi(N)=N-(p+q)+1$ we recover $p+q=N-\varphi(N)+1=:s$. Then $p,q$ are the roots of
$t^{2}-st+N=0$, i.e.\ $p,q=\tfrac12\bigl(s\pm\sqrt{s^{2}-4N}\bigr)$ --- a constant-time
computation. \end{proof}

Thus the entire security of RSA rests on the presumed hardness of the
\textbf{integer-factorisation problem} (and the slightly weaker ``RSA problem'':
computing $e$-th roots modulo $N$). No efficient \emph{classical} factoring algorithm
is known; the best general method, the General Number Field Sieve, runs in
sub-exponential but super-polynomial time $\exp\!\big(O(\log N)^{1/3}(\log\log N)^{2/3}\big)$.

\begin{remark}[Textbook vs.\ real RSA]
``Textbook RSA'' as above is deterministic and therefore \emph{insecure} in practice
(equal messages give equal ciphertexts, small $m$ with small $e$ may satisfy $m^e<N$ so
$c=m^e$ exactly, etc.). Real deployments use randomised padding (RSA-OAEP for
encryption, RSA-PSS for signatures) and $e=65537$. These do not change the number
theory above; they protect against structural attacks.
\end{remark}

% ============================================================================
\section{Shor's algorithm}
% ============================================================================

Shor's algorithm (1994) factors $N$ in time polynomial in $\log N$ on a quantum
computer, breaking RSA asymptotically. Its heart is a purely classical reduction
(factoring $\to$ order finding) wrapped around one quantum subroutine (order finding
via phase estimation).

\subsection{Step 1 --- Reduce factoring to order finding}

\begin{definition}[Multiplicative order]
For $a$ with $\gcd(a,N)=1$, the \emph{order} $r=\ord_N(a)$ is the smallest positive
integer with $a^{r}\equiv 1 \pmod N$.
\end{definition}

\begin{theorem}[Order finding yields factors]\label{thm:order-factor}
Let $N$ be odd and composite, $a$ coprime to $N$ with order $r$. If
\begin{enumerate}[label=(\roman*),leftmargin=2.4em,itemsep=0pt,topsep=2pt]
  \item $r$ is \textbf{even}, and
  \item $a^{r/2}\not\equiv -1 \pmod N$ \ (note $a^{r/2}\not\equiv +1$ automatically, by
        minimality of $r$),
\end{enumerate}
then $\gcd\!\big(a^{r/2}-1,\,N\big)$ and $\gcd\!\big(a^{r/2}+1,\,N\big)$ are
\textbf{non-trivial factors} of $N$.
\end{theorem}

\begin{proof}
Set $x=a^{r/2}\bmod N$. Then $x^{2}=a^{r}\equiv 1$, so
\[
  N \mid x^{2}-1 = (x-1)(x+1).
\]
If $N$ divided $x-1$ we would have $x\equiv 1$, contradicting minimality of $r$; if $N$
divided $x+1$ we would have $x\equiv -1$, excluded by (ii). Hence $N$ divides the product
$(x-1)(x+1)$ but neither factor alone, so its prime factors are split between the two.
Therefore $\gcd(x-1,N)$ is a proper divisor $>1$ of $N$ (and so is $\gcd(x+1,N)$).
\end{proof}

\begin{lemma}[The gamble usually pays off]\label{lem:prob}
If $N$ has $k$ distinct odd prime factors and $a$ is drawn uniformly from $\Zn$, then
conditions (i)--(ii) of Theorem~\ref{thm:order-factor} both hold with probability at
least $1-2^{-k}\ge \tfrac12$.
\end{lemma}

So on average two or three random choices of $a$ suffice; if $a$ fails, pick another.

\subsection{Step 2 --- Compute the order on a quantum computer}

The only hard step is finding $r=\ord_N(a)$. Classically this is as hard as factoring;
quantumly it is \emph{easy}, via Quantum Phase Estimation (QPE) of the operator
\[
  U\ket{y} = \ket{ay \bmod N}, \qquad y\in\{0,1,\dots,N-1\}.
\]
$U$ is unitary (multiplication by $a$ permutes the residues coprime to $N$).

\begin{proposition}[Eigenstructure of $U$]\label{prop:eig}
For each $s\in\{0,1,\dots,r-1\}$ the state
\[
  \ket{u_s} = \frac{1}{\sqrt r}\sum_{k=0}^{r-1} e^{-2\pi i\, s k / r}\,\ket{a^{k}\bmod N}
\]
is an eigenvector of $U$ with eigenvalue $e^{2\pi i\, s/r}$:
\[
  U\ket{u_s} = e^{2\pi i\, s/r}\ket{u_s}.
\]
Moreover
\[
  \frac{1}{\sqrt r}\sum_{s=0}^{r-1}\ket{u_s} = \ket{1}.
\]
\end{proposition}

\begin{proof}[Idea]
Applying $U$ shifts $a^{k}\mapsto a^{k+1}$ in the sum, which is equivalent to multiplying
the whole state by $e^{2\pi i s/r}$ after re-indexing (the phases were chosen precisely so
that this works). Summing $\ket{u_s}$ over $s$, the phase factors
$\sum_s e^{-2\pi i s k/r}$ vanish for $k\not\equiv 0$ and reinforce for $k=0$, leaving only
$\ket{a^{0}}=\ket{1}$. \end{proof}

The eigenphase of each $\ket{u_s}$ is $\phi_s = s/r$ --- it \emph{contains the order}. Two
facts make this usable:
\begin{itemize}[leftmargin=1.6em,itemsep=1pt,topsep=2pt]
  \item We cannot prepare a single $\ket{u_s}$ (we don't know $r$), but we \emph{can}
        prepare $\ket{1}=\tfrac{1}{\sqrt r}\sum_s\ket{u_s}$ trivially.
  \item Running QPE on $U$ with input $\ket{1}$ therefore measures a phase
        $\phi\approx s/r$ for a \textbf{uniformly random} $s\in\{0,\dots,r-1\}$.
\end{itemize}

\paragraph{The QPE circuit.} With $t$ \emph{counting} qubits (initialised to
$\ket{+}^{\otimes t}$) and a work register holding $\ket{1}$, apply controlled
$U^{2^{0}},U^{2^{1}},\dots,U^{2^{t-1}}$, then the inverse Quantum Fourier Transform on the
counting register, then measure it. The read-out integer $y$ estimates
$\phi\approx y/2^{t}\approx s/r$. Using $t\approx 2\log_2 N + 1$ counting qubits makes the
estimate sharp enough that continued fractions recover $r$ with constant probability.

\subsection{Step 3 --- Recover $r$ by continued fractions}

QPE returns a binary fraction $\phi=y/2^{t}$ close to the rational $s/r$. We need $s/r$ in
lowest terms.

\begin{theorem}[Continued-fraction recovery]\label{thm:cf}
If a rational $s/r$ with $0\le s<r\le N$ is approximated by $\phi$ to within
$|\,\phi - s/r\,| \le \dfrac{1}{2r^{2}}$, then $s/r$ is one of the \emph{convergents} of
the continued-fraction expansion of $\phi$, and is found in $O(\log N)$ steps.
\end{theorem}

In code this is a single call: \verb|Fraction(phi).limit_denominator(N)| returns the best
rational approximation to $\phi$ with denominator $\le N$ --- exactly the convergent
guaranteed above. The denominator of that fraction is the candidate for $r$.

\begin{remark}[Why the algorithm is probabilistic]
Three things can go wrong on a given run: the random $s$ may be $0$ (useless), $s$ and $r$
may share a factor (we then recover a \emph{divisor} of $r$, not $r$ itself), or the
recovered $r$ may be odd / give $a^{r/2}\equiv-1$. Each failure is detected cheaply
($a^{r}\stackrel{?}{\equiv}1$, and the gcd test), and we simply repeat. A constant number
of repetitions gives any desired success probability.
\end{remark}

\subsection{The complete recipe}

\begin{recipebox}{Recipe S1 --- Shor's factoring algorithm for $N$}
\begin{enumerate}[leftmargin=2.2em,itemsep=1pt,topsep=3pt]
  \item \textbf{Classical shortcuts.} If $N$ is even, output $2$. If $N=c^{b}$ is a perfect
        power (checkable classically), output $c$.
  \item Pick a random $a\in\{2,\dots,N-1\}$. Compute $g=\gcd(a,N)$.
        If $g>1$, output $g$ (a lucky factor).
  \item \textbf{Quantum order finding.} Use QPE on $U\ket{y}=\ket{ay\bmod N}$ with input
        $\ket{1}$ to measure $\phi\approx s/r$.
  \item \textbf{Continued fractions.} Set $r=\operatorname{denominator}$ of
        $\texttt{limit\_denominator}(\phi,\,N)$ (Theorem~\ref{thm:cf}).
        Verify $a^{r}\equiv 1 \pmod N$; if not, go to step 2.
  \item \textbf{Extract factors.} If $r$ is even and $a^{r/2}\not\equiv -1 \pmod N$, output
        $\gcd(a^{r/2}\pm1,\,N)$ (Theorem~\ref{thm:order-factor}). Otherwise go to step 2.
\end{enumerate}
\end{recipebox}

\subsection{Worked quantum examples}

\begin{example}[Factoring $N=15$ with $a=7$]\label{ex:15}
The order of $7$ modulo $15$ is $r=4$ (since $7^{1}=7,\,7^{2}=49\equiv4,\,7^{3}\equiv13,\,
7^{4}\equiv1$). QPE with $t=8$ counting qubits returns one of the phases
$\phi\in\{0,\tfrac14,\tfrac12,\tfrac34\}$, i.e.\ register values $\{0,64,128,192\}$ out of
$256$. Taking $\phi=\tfrac34$: \verb|Fraction(0.75).limit_denominator(15)|$=3/4$, so $r=4$.
Then
\[
  a^{r/2}=7^{2}=49\equiv 4 \pmod{15},\quad
  \gcd(4-1,15)=3,\quad \gcd(4+1,15)=5.
\]
We recover $15 = 3\times 5$. (The outcomes $\phi=0$ and $\phi=\tfrac12$ fail and are
retried --- this is exactly the histogram you see in the notebook.)
\end{example}

\begin{example}[Factoring $N=21$ with $a=2$]\label{ex:21}
Powers of $2$ mod $21$: $2,4,8,16,11,1$, so $r=6$ (even). Then
$a^{r/2}=2^{3}=8$, and
\[
  \gcd(8-1,21)=\gcd(7,21)=7,\qquad \gcd(8+1,21)=\gcd(9,21)=3,
\]
giving $21 = 3\times 7$.
\end{example}

\begin{example}[Breaking the toy RSA key of Example~\ref{ex:rsa} by factoring]
Suppose an attacker sees only the public key $(N,e)=(55,3)$ and the ciphertext $c=13$.
Running Recipe~S1 on $N=55$ recovers $p=5,q=11$. Then
$\varphi(N)=40$, $d\equiv 3^{-1}\equiv 27 \pmod{40}$, and
$m = c^{d}\bmod N = 13^{27}\bmod 55 = 7$. The private key --- and the message --- are
recovered purely from public data. This is the attack \texttt{16\_Shor\_Period\_Finding\_to\_RSA.ipynb} performs on a
simulator in its ``breaking RSA'' section.
\end{example}

\subsection{The bottleneck: modular exponentiation}

The controlled-$U^{2^{j}}$ gates dominate the cost. We never build $U^{2^{j}}$ by repeating
$U$ a full $2^{j}$ times (that is exponential); instead
$a^{2^{j}}\bmod N$ is precomputed \emph{classically} by repeated squaring in $O(j)$
multiplications, and the modular multiplier is realised as a reversible circuit. Beauregard's
construction factors an $n$-bit $N$ with only $2n+3$ qubits and $O(n^{3})$ gates.

\begin{theorem}[Complexity of Shor's algorithm]
Order finding, and hence factoring $N$, succeeds with constant probability using
$O(\log N)$ qubits and $\widetilde{O}\big((\log N)^{2}\big)$ quantum gates per run (dominated
by modular exponentiation), plus $O((\log N)^{3})$ classical post-processing --- polynomial
in the input size $\log N$.
\end{theorem}

\subsection{Reality check and post-quantum cryptography}

Although Shor's algorithm is polynomial, the \emph{constants} and the cost of quantum error
correction are enormous. The widely cited estimate of Gidney and Eker\aa\ (2019) puts the
cost of factoring a 2048-bit RSA modulus at roughly $20$ million \emph{noisy physical} qubits
running for about $8$ hours, corresponding to a few thousand logical (error-corrected) qubits.
Today's devices --- including IQM \emph{Garnet} (20 qubits) and \emph{Emerald} (54 qubits), and
the largest superconducting chips (a few hundred physical qubits) --- are many orders of
magnitude short of this. Hence the companion notebooks factor $15$ and $21$, not real keys.

The threat is nevertheless taken seriously (``harvest now, decrypt later''), which is why NIST
has standardised \textbf{post-quantum} schemes whose security does \emph{not} reduce to
factoring or discrete logs --- notably the lattice-based ML-KEM (Kyber) and ML-DSA (Dilithium).

% ============================================================================
\section{From the lesson to the notebook}
% ============================================================================

The companion notebook \texttt{16\_Shor\_Period\_Finding\_to\_RSA.ipynb} realises this lesson cell by cell:
\begin{itemize}[leftmargin=1.6em,itemsep=1pt,topsep=2pt]
  \item \S2 $\leftrightarrow$ \texttt{factors\_from\_order}: Theorem~\ref{thm:order-factor}
        in code.
  \item \S4.2 $\leftrightarrow$ \texttt{build\_order\_finding\_circuit},
        \texttt{c\_amod15}, \texttt{qft\_dagger}: the QPE circuit of
        Proposition~\ref{prop:eig}.
  \item \S4.3 $\leftrightarrow$ \texttt{best\_rational} / \texttt{limit\_denominator}:
        continued-fraction recovery (Theorem~\ref{thm:cf}).
  \item Recipe~S1 $\leftrightarrow$ \texttt{shor\_factor(N)}: the full loop.
  \item Example~\ref{ex:rsa} and the ``breaking RSA'' cell $\leftrightarrow$
        Theorem~\ref{thm:rsa} and Proposition~\ref{prop:phi-fact}.
\end{itemize}
Run it on the Aer simulator (\texttt{IBM/16\_Shor\_Period\_Finding\_to\_RSA.ipynb}) or on IQM Garnet/Emerald
(\texttt{IQM/16\_Shor\_Period\_Finding\_to\_RSA.ipynb}) by changing a single \texttt{TARGET} line.

% ============================================================================
\section*{Exercises}
% ============================================================================
\addcontentsline{toc}{section}{Exercises}
\begin{enumerate}[leftmargin=2em,itemsep=2pt]
  \item Compute $\varphi(561)$ and verify that $561=3\cdot 11\cdot 17$ is a Carmichael
        number ($a^{560}\equiv1$ for all $a$ coprime to $561$, yet $561$ is composite).
  \item For $p=7,q=13$ build an RSA key with $e=5$; find $d$ and encrypt/decrypt $m=9$.
  \item Prove that if $e$ is chosen with $\gcd(e,\varphi(N))\neq 1$ then $d$ need not exist,
        and explain why Recipe~R1 step~3 forbids it.
  \item In Example~\ref{ex:15}, which measured phases $\phi$ \emph{fail} to produce a factor,
        and why? Relate your answer to the Remark on probabilistic failure.
  \item Factor $N=35$ by hand using $a=3$ ($r=12$): show $\gcd(3^{6}\pm1,35)$ gives
        $\{5,7\}$.
\end{enumerate}

% ============================================================================
\begin{thebibliography}{9}
\addcontentsline{toc}{section}{References}
\bibitem{rsa} R.~Rivest, A.~Shamir, L.~Adleman,
  \emph{A Method for Obtaining Digital Signatures and Public-Key Cryptosystems},
  Comm.\ ACM \textbf{21} (1978) 120--126.
\bibitem{shor} P.~W.~Shor,
  \emph{Polynomial-Time Algorithms for Prime Factorization and Discrete Logarithms on a
  Quantum Computer}, SIAM J.\ Comput.\ \textbf{26} (1997) 1484;
  \href{https://arxiv.org/abs/quant-ph/9508027}{arXiv:quant-ph/9508027}.
\bibitem{nc} M.~A.~Nielsen and I.~L.~Chuang,
  \emph{Quantum Computation and Quantum Information}, Cambridge Univ.\ Press (2000), Ch.~5.
\bibitem{beauregard} S.~Beauregard,
  \emph{Circuit for Shor's algorithm using $2n+3$ qubits},
  \href{https://arxiv.org/abs/quant-ph/0205095}{arXiv:quant-ph/0205095}.
\bibitem{gidney} C.~Gidney and M.~Eker\aa,
  \emph{How to factor 2048 bit RSA integers in 8 hours using 20 million noisy qubits},
  Quantum \textbf{5} (2021) 433; \href{https://arxiv.org/abs/1905.09749}{arXiv:1905.09749}.
\end{thebibliography}

\end{document}
